\documentclass{zhvt-classic}

% 使用本地源流明体（GenRyuMin2 TW）按文件路径加载，不依赖系统字体名
% 版框（内框／界栏／书口侧栏／鱼尾）改用与版心印章同一印泥棕色：
%   印章图片 yshs.png 实色 = RGB(135,68,52) = HTML 874434
\usepackage{xcolor}
\usepackage{tikz}
% sealbrown 双注册：xcolor 供 TikZ/pgf，L3 供 cls 版框
\definecolor{sealbrown}{HTML}{874434}
\ExplSyntaxOn
\color_set:nnn { sealbrown } { HTML } { 874434 }
\ExplSyntaxOff

\zhvtset{
  grid color       = sealbrown,
  main font        = GenRyuMin2TW-SB,
  main font path   = C:/Users/Administrator/texmf/fonts/opentype/public/GenRyuMin2TW/,
jiazhu font      = GenRyuMin2TW-M,
  jiazhu font path = C:/Users/Administrator/texmf/fonts/opentype/public/GenRyuMin2TW/,
  % 页码框容纳 4 字：正常页「九十七」+前置空 1 格=4 行恰好；
  % 100+ 页「一百零一」共 5 行超出，末字「一/二/三」被框裁掉（2026-10-04 实测），故放宽到 5
  shukou page number chars = 5,
  % 书口书名框高从默认 6 字减到 4 字，把省出的 2 字高拨给 chapter 框，
  % 使 6 字书名+卷名放进 9 字高的 chapter 框后不再压线（2026-10-05）
  shukou title chars = 4,
}

\title{史记}[孟子荀卿列传]
\author{汉·司马迁}
\maker{闲叟漫记·重製}


% ===== 索隐述赞样式（2026-11-08 四次定稿 · 愚大人亲手改定参数）=====
% ⚠ 竖排前提：自定义字体家族必须带 RawFeature={vertical:+vert}；
%   simhei.ttf 加 +vhal 会整体加载失败静默回退，只写 +vert；simfang 同理。
\newCJKfontfamily{\shuzanfont}[Path=C:/Windows/Fonts/,RawFeature={vertical:+vert}]{simhei.ttf}
\newCJKfontfamily{\shuzanbody}[Path=C:/Windows/Fonts/,RawFeature={vertical:+vert}]{simfang.ttf}
% 题头：152x30pt 圆角 6pt 棕底白字；baseline (0,2) 使整框落入界行格子不压版框线；
% 白字节点 x 不补偿、y +8.2pt（像素级墨迹实测 dx/dy≈0，L/R 留白约 20pt）。
\newcommand{\shuzanhead}[1]{%
  \begin{tikzpicture}[baseline={(0,2)}]
    \fill[sealbrown, rounded corners=6pt]
      (-76pt,-15pt) rectangle (76pt,15pt);
    \node[text=white, inner sep=0pt, anchor=center] at (0pt,8.2pt)
      {\shuzanfont\fontsize{23}{23}\selectfont #1};
  \end{tikzpicture}%
}
% ★ 修正书口卷名位置：chapter 框起始位置改回 cls 默认值，避免全局影响。
%   只在每个卷首页通过 \juan 命令里的空命令微调该页位置（安全，不会跨页污染）。
\ExplSyntaxOn
\AtBeginDocument{
  \dim_set:Nn \l__zhvt_shukou_chapter_title_start_dim
    { \l__zhvt_shukou_title_height_dim + \l__zhvt_font_size_dim * 2 }
  \dim_set:Nn \l__zhvt_shukou_page_number_start_dim
    { \l__zhvt_shukou_chapter_title_start_dim + \l__zhvt_shukou_chapter_title_height_dim }
}
\ExplSyntaxOff

\begin{document}

% ★★ 版心（书口）排布覆写：把书名放到鱼尾之下（2026-10-04）
%
%   cls 默认版心自上而下：书名（title 静态框，在鱼尾之上）→ 鱼尾 → 篇名（chapter 动态框）→ 页码
%   本模板改为：          鱼尾（位置不动）→ 书名 → 篇名 → 页码分割线 → 页码
%
%   手法：把 cls 放在 title 静态框里的书名，改放进 chapter 动态框（鱼尾之下）；
%         title 静态框留空。书名取 \l__zhvt_title_str，即 \title{...} 的第一参数，
%         所以换书时只改 \title，不必在此重复写书名。
%
%   ⚠ 三条必须遵守的写法（2026-10-04 实测）：
%     ① 汉字必须用 \hbox{...} 包住。flowfram 会把框内容装进一个宽 14.4mm（= 一个 line_sep
%        = 书口栏宽）的纵向 minipage 并按该宽度折行，而整页靠 PDF /Rotate 90 显示、
%        框内文字本身不变向 —— 不包 \hbox 时「围炉夜话」4 字会被逐字折成 4 行，
%        在版面上就成了横排 4 字并横向溢出版框。
%     ② 竖排下**一个 \hbox 就是一整列**。书名与篇名要同处一列，必须放进同一个 \hbox；
%        分成两个 \hbox 时，后一个会被当成新的一列、甩到隔壁栏去
%        （实测「序」跑到书口栏外侧 14.4mm 的正文列上）。中间的 \quad 在框内是横向空白，
%        显示出来正好是书名与篇名之间的纵向间隔。
%     ③ ★ 内容盒要写成 **与框等宽的定宽盒** `\hbox to \l__zhvt_shukou_chapter_title_height_dim{…\hfil}`。
%        实测：chapter 框是把内容盒沿纵向**居中**摆放的，内容宽度一变，书名就跟着漂：
%        自然宽度的 \hbox（书名 38.4mm）落在 x=157.2mm，加篇名后（57.6mm）落在 x=152.4mm，
%        目录页（无篇名）与正文页的书名就差半格。做成等宽盒后，无论有没有篇名，
%        内容固定从框顶起排，实测两页均落在 x=142.8mm，不再漂。
\ExplSyntaxOn
% 把书口 chapter 框（书名+篇名）起始位置下移 1 字高，避免 4 字卷名顶到鱼尾
% ⚠ cls 的 \__zhvt_calculate_dimensions: 在 \AtBeginDocument 里执行，会重算该 dim；
%   必须把覆写挂到 begin document 之后，否则会被二次覆盖（2026-10-05 实测）。
% (旧方案已移到导言区，直接修改 chapter 框起始位置)
% 卷标题命令 #1=书口篇名(鱼尾) #2=正文卷名（含卷序，卷名内空格用 ~）
%
% ★★ 目录页两项修订（2026-10-05，愚大人指出）★★
%   ① 页码起点统一对齐：
%      原状 —— 目录项写成 \par\hspace{4em}卷名\hspace{4em}页码，页码起点 = 卷名实宽 + 4em。
%      卷名实宽随字数变：「卷一~梁惠王上」6 字 +~、「卷七~离娄上」5 字 +~，
%      差 1 个字位 → 卷七~卷十 4 条的页码整体比其余 10 条高 27.21pt（一个字），看着上下不齐。
%      改法 —— 卷名补全角空格（宽 = 1 字 = 27.215pt）补足到 **7 字符位**再写进目录，
%      14 条页码起点遂全部落在同一条竖线上（起点对齐；末字随页码字数自然参差）。
%      ⚠ 只对目录里的副本 \l_tmpc_tl 补，正文卷名仍用 \l_tmpa_tl 原样 —— 正文像素不受影响。
%   ② 卷名整体上移半字位：前导 \hspace 由 4em 改 3.5em。
%      实测原状目录块在版框内偏下（天头留白 139.6pt / 地脚留白 110.7pt），
%      上移 0.5em = 13.6pt 后上下基本均分，视觉居中。
\NewDocumentCommand { \juan } { mm } {
	\clearpage
	\tl_set:Nn \l__zhvt_chapter_title_tl { #1 }
	\tl_set:Nn \l_tmpa_tl { #2 }
	\tl_set_eq:NN \l_tmpc_tl \l_tmpa_tl
	\int_step_inline:nn { 7 - \str_count:n { #2 } }
	  { \tl_put_right:Nn \l_tmpc_tl { 　 } }
	\tl_set:Nx \l_tmpb_tl { \zhnum { page } }
	\tl_put_left:Nn \l_tmpb_tl { \hspace { 4em } }
	\tl_put_left:Nn \l_tmpb_tl { \l_tmpc_tl }
	\tl_put_left:Nn \l_tmpb_tl { \par \hspace { 3.5em } }
	\addtocontents { toc } { \l_tmpb_tl }
	\skip_horizontal:n { \l__zhvt_font_size_dim * 4 }
	\l_tmpa_tl
	\par
	\thispagestyle { empty }
}
\cs_set_protected:Npn \__zhvt_set_shukou_content: {
	% ★ 鱼尾上方留白（title 静态框）嵌入闲叟钤印：
	%   整页 /Rotate 90，故此处把图用 \rotatebox{90} 转正为竖章；
	%   尺寸卡在版心条幅内（书内约 7mm 宽 × 33.8mm 高，留余量不压鱼尾）。
	\setstaticcontents * { title } {
		\vbox_to_ht:nn { \l__zhvt_baseline_skip_dim } {
			\vskip 14.2pt
			\hbox_to_wd:nn { \l__zhvt_shukou_title_height_dim } {
				\hfil
				\rotatebox [ origin = c ] { 90 } { \includegraphics [ width = 7mm , height = 33.8mm ] { yshs.png } }
				\hfil
			}
		}
	}
% ★ 版心（书口）统一改法（2026-11-16 愚大人定）：18pt＋框内居中＋书口文字同版框色
%   ① 原 27.2pt 时 7 字长篇名（如「廉颇蔺相如列传」）末字溢出 chapter 框、压住页码；
%      改 \fontsize{18pt}{18pt}，前后 \hfil 令字串在框内居中（不随篇名字数漂移）。
%   ② \vbox{\vskip 8.2pt} 定量下移，抵消缩字号后栏宽方向 8.2pt 上偏（像素实测 y 中心 543.3 = 原版）。
%      ⚠ \raisebox / \vspace* 在 flowfram 框内无效；换字号须同步换补偿：16pt→10.1、19pt→7.4。
%   ③ chapter 与 page~number 均置 \color{sealbrown}（HTML 874434），与版框／界行同色；正文仍黑。
%   ⚠ 页码字位、钤印、鱼尾、界行、正文一律未动；五部合订本同此口径。
	\setdynamiccontents * { chapter }
	{ \color{sealbrown} \vbox { \vskip 8.2pt \hbox to \l__zhvt_shukou_chapter_title_height_dim
		{ \hfil \fontsize{18pt}{18pt}\selectfont \l__zhvt_title_str \quad \l__zhvt_chapter_title_tl \hfil } } }
	\setdynamiccontents * { page~number } { \color{sealbrown} \quad \zhnum { page } }
}
\__zhvt_set_shukou_content:
\ExplSyntaxOff

\maketitle

% 单篇，无多卷目录，不渲染目录页

\juan{孟子荀卿列传}{卷七十四~孟子荀卿列传第十四}
\hspace*{2em}太史公曰。余读孟子书。至梁惠王问﹁何以利吾国﹂。未尝不废书而叹也。曰。嗟乎。利诚乱之始也。夫子罕言利者。常防其原也。故曰﹁放於利而行。多怨﹂。自天子至於庶人。好利之弊何以异哉。\\
\hspace*{2em}孟轲。驺人也。受业子思之门人。道既通。游事齐宣王。宣王不能用。適梁。梁惠王不果所言。则见以为迂远而阔於事情。当是之时。秦用商君。富国彊兵。楚。魏用吴起。战胜弱敌。齐威王。宣王用孙子。田忌之徒。而诸侯东面朝齐。天下方务於合从连衡。以攻伐为贤。而孟轲乃述唐。虞。三代之德。是以所如者不合。退而与万章之徒序诗书。述仲尼之意。作孟子七篇。其後有驺子之属。\\
\hspace*{2em}齐有三驺子。其前驺忌。以鼓琴干威王。因及国政。封为成侯而受相印。先孟子。\\
\hspace*{2em}其次驺衍。後孟子。驺衍睹有国者益淫侈。不能尚德。若大雅整之於身。施及黎庶矣。乃深观阴阳消息而作怪迂之变。终始。大圣之篇十馀万言。其语闳大不经。必先验小物。推而大之。至於无垠。先序今以上至黄帝。学者所共术。大并世盛衰。因载其禨祥度制。推而远之。至天地未生。窈冥不可考而原也。先列中国名山大川。通谷禽兽。水土所殖。物类所珍。因而推之。及海外人之所不能睹。称引天地剖判以来。五德转移。治各有宜。而符应若兹。以为儒者所谓中国者。於天下乃八十一分居其一分耳。中国名曰赤县神州。赤县神州内自有九州。禹之序九州是也。不得为州数。中国外如赤县神州者九。乃所谓九州也。於是有裨海环之。人民禽兽莫能相通者。如一区中者。乃为一州。如此者九。乃有大瀛海环其外。天地之际焉。其术皆此类也。然要其归。必止乎仁义节俭。君臣上下六亲之施。始也滥耳。王公大人初见其术。惧然顾化。其後不能行之。\\
\hspace*{2em}是以驺子重於齐。適梁。惠王郊迎。执宾主之礼。適赵。平原君侧行撇席。如燕。昭王拥彗先驱。请列弟子之座而受业。筑碣石宫。身亲往师之。作主运。其游诸侯见尊礼如此。岂与仲尼菜色陈蔡。孟轲困於齐梁同乎哉。故武王以仁义伐纣而王。伯夷饿不食周粟。卫灵公问陈。而孔子不答。梁惠王谋欲攻赵。孟轲称大王去邠。此岂有意阿世俗苟合而已哉。持方枘欲内圜凿。其能入乎。或曰。伊尹负鼎而勉汤以王。百里奚饭牛车下而缪公用霸。作先合。然後引之大道。驺衍其言虽不轨。傥亦有牛鼎之意乎。\\
\hspace*{2em}自驺衍与齐之稷下先生。如淳于髡。慎到。环渊。接子。田骈。驺奭之徒。各著书言治乱之事。以干世主。岂可胜道哉。\\
\hspace*{2em}淳于髡。齐人也。博闻彊记。学无所主。其谏说。慕晏婴之为人也。然而承意观色为务。客有见髡於梁惠王。惠王屏左右。独坐而再见之。终无言也。惠王怪之。以让客曰。﹁子之称淳于先生。管。晏不及。及见寡人。寡人未有得也。岂寡人不足为言邪。何故哉。﹂客以谓髡。髡曰。﹁固也。吾前见王。王志在驱逐。後复见王。王志在音声。吾是以默然。﹂客具以报王。王大骇。曰。﹁嗟乎。淳于先生诚圣人也。前淳于先生之来。人有献善马者。寡人未及视。会先生至。後先生之来。人有献讴者。未及试。亦会先生来。寡人虽屏人。然私心在彼。有之。﹂後淳于髡见。壹语连三日三夜无倦。惠王欲以卿相位待之。髡因谢去。於是送以安车驾驷。束帛加璧。黄金百镒。终身不仕。\\
\hspace*{2em}慎到。赵人。田骈。接子。齐人。环渊。楚人。皆学黄老道德之术。因发明序其指意。故慎到著十二论。环渊著上下篇。而田骈。接子皆有所论焉。\\
\hspace*{2em}驺奭者。齐诸驺子。亦颇采驺衍之术以纪文。\\
\hspace*{2em}於是齐王嘉之。自如淳于髡以下。皆命曰列大夫。为开第康庄之衢。高门大屋。尊宠之。览天下诸侯宾客。言齐能致天下贤士也。\\
\hspace*{2em}荀卿。赵人。年五十始来游学於齐。驺衍之术迂大而闳辩。奭也文具难施。淳于髡久与处。时有得善言。故齐人颂曰。﹁谈天衍。雕龙奭。炙毂过髡。﹂田骈之属皆已死齐襄王时。而荀卿最为老师。齐尚脩列大夫之缺。而荀卿三为祭酒焉。齐人或谗荀卿。荀卿乃適楚。而春申君以为兰陵令。春申君死而荀卿废。因家兰陵。李斯尝为弟子。已而相秦。荀卿嫉浊世之政。亡国乱君相属。不遂大道而营於巫祝。信禨祥。鄙儒小拘。如庄周等又猾稽乱俗。於是推儒。墨。道德之行事兴坏。序列著数万言而卒。因葬兰陵。\\
\hspace*{2em}而赵亦有公孙龙为坚白同异之辩。剧子之言。魏有李悝。尽地力之教。楚有尸子。长卢。阿之吁子焉。自如孟子至于吁子。世多有其书。故不论其传云。\\
\hspace*{2em}盖墨翟。宋之大夫。善守御。为节用。或曰并孔子时。或曰在其後。\\

% ===== 索隐述赞（司马贞《史记索隐》）· 统一样式 =====
% 题头「索隐述赞曰」五字并入反白框；正文每四句一列，定宽盒 \hfil 居列中；
% \fontsize{24}{41} 列距对齐全书界行（121+41k）；\clearpage 使述赞独占一页。
\clearpage
\shuzanhead{索隐述赞曰}\linebreak
{\shuzanbody\fontsize{24}{41}\selectfont
\hbox to \hsize{\hfil 六国之末。战胜相雄。轲游齐。魏。\hfil}\linebreak
\hbox to \hsize{\hfil 其说不通。退而著述。称吾道穷。兰陵事楚。\hfil}\linebreak
\hbox to \hsize{\hfil 驺衍谈空。康庄虽列。莫见收功。\hfil}%
\par}
\end{document}
